% !TEX root = ../../main.tex

\renewcommand{\thesection}{2.3}

\section{Laplace 定理}\label{sec:2.3}

\begin{ex}{3}\label{ex:2.3.3}
    设 \(D=
    \begin{vmatrix}
        a_1 & a_2 & \cdots & a_n\\
        a_2 & 1 & &\\
        \vdots & & \ddots &\\
        a_n & & & 1
    \end{vmatrix}
    \), 求 \(D\) 的第一列元素的代数余子式之和.
\end{ex}

\begin{solution}
    设 \(D=\vsmallmatrix{\alpha_1 & \alpha_2 & \cdots & \alpha_n},D'=\vsmallmatrix{1 & \alpha_2 & \cdots & \alpha_n}\),
    由按列展开定理, \(D'\) 即为所求.

    由行列式的性质,
    \begin{align*}
        D'&=\vsmallmatrix{1 & \alpha_2 & \cdots & \alpha_n}\\
        &=\vsmallmatrix{1-\sum_{i=2}^{n}\alpha_i & \alpha_2 & \cdots & \alpha_n}\\
        &=
        \begin{vmatrix}
            1-\sum\limits_{i=2}^{n}a_i & a_2 & \cdots & a_n\\
            & 1 & &\\
            & & \ddots &\\
            & & & 1
        \end{vmatrix}\\
        &=1-\sum_{i=2}^{n}a_i.\qedhere
    \end{align*}
\end{solution}

\begin{ex}{5}\label{ex:2.3.5}
    计算 \(2n\) 阶行列式 \(
    \begin{vmatrix}
        a & & & & & b\\
        & \ddots & & & \iddots &\\
        & & a & b & &\\
        & & b & a & &\\
        & \iddots & & & \ddots &\\
        b & & & & & a
    \end{vmatrix}
    \).
\end{ex}

\begin{solution}
    设行列式数列 \(\{D_n\}\) 满足 \(D_1=
    \begin{vmatrix}
        a & b\\
        b & a
    \end{vmatrix}
    \), \(D_{n+1}=
    \begin{vmatrix}
        a & & b\\
        & D_n &\\
        b & & a
    \end{vmatrix}
    \) (\(n\in\N^*\)).

    由 Laplace 定理, 有
    \[
    D_{n+1}=
    \begin{vmatrix}
        a & b\\
        b & a
    \end{vmatrix}
    \cdot(-1)^{1+2n+1+2n}D_n=D_1D_n.
    \]
    故有通项公式
    \[
    D_n=D_1^n=(a^2-b^2)^n.\qedhere
    \]
\end{solution}

\begin{ex}{7}\label{2.3.7}
    设 \(a_1,a_2,\cdots,a_n\) 互异, 求多项式 \(f(x)=
    \begin{vmatrix}
        1 & 1 & \cdots & 1\\
        x & a_1 & \cdots & a_n\\
        \vdots & \vdots & & \vdots\\
        x^n & a_1^n & \cdots & a_n^n
    \end{vmatrix}
    \) 的根.
\end{ex}

\begin{solution}
    该行列式为 Vandermonde 行列式, 故
    \[
    f(x)=\prod_{1\leq i<j\leq n}(a_j-a_i)\cdot\prod_{i=1}^n(x-a_i),
    \]
    故 \(f(x)\) 的根有 \(a_1,a_2,\cdots,a_n\).
\end{solution}

\begin{ex}{8}\label{2.3.8}
    求下面 \(n\) 阶行列式 \(\det(a_{ij})\).
    \[
    a_{ij}=a_ib_j+\delta_{ij}c_i.
    \]
\end{ex}

\begin{solution}
    设
    \[    
    \begin{vmatrix}
        c_1 & 0 & \cdots & 0\\
        0 & c_2 & \cdots & 0\\
        \vdots & \vdots & \ddots & \vdots\\
        0 & 0 & \cdots & c_n
    \end{vmatrix}
    =\vsmallmatrix{\gamma_1 & \gamma_2 & \cdots & \gamma_n},
    \]
    \begin{align*}
        \det(a_{ij})&=
        \begin{vmatrix}
            a_1b_1+c_1 & a_1b_2 & \cdots & a_1b_n\\
            a_2b_1 & a_2b_2+c_2 & \cdots & a_2b_n\\
            \vdots & \vdots & \ddots & \vdots\\
            a_nb_1 & a_nb_2 & \cdots & a_nb_n+c_n
        \end{vmatrix}\\
        &=\vsmallmatrix{b_1\alpha+\gamma_1 & b_2\alpha+\gamma_2 & \cdots & b_n\alpha+\gamma_n}.
    \end{align*}
    则有
    \begin{align*}
        \det(a_{ij})
        &=\vsmallmatrix{b_1\alpha+\gamma_1 & b_2\alpha+\gamma_2 & \cdots & b_n\alpha+\gamma_n}\\
        &=\sum_{S\in\{1,2,\cdots,n\}}\vsmallmatrix{\beta_S(1) & \beta_S(2) & \cdots & \beta_S(n)},
    \end{align*}
    其中
    \[
    \beta_S(t)=
    \begin{cases}
    b_t\alpha,& t\in S,\\
    \gamma_t,& t\notin S.
    \end{cases}
    \]

    当 \(|S|\geq2\), 设 \(i,j\in S\) 满足 \(i\neq j\), 则
    \begin{align*}
        \vsmallmatrix{\beta_S(1) & \beta_S(2) & \cdots & \beta_S(n)}
        &=\vsmallmatrix{\cdots & b_i\alpha & \cdots & b_j\alpha & \cdots}\\
        &=b_ib_j\vsmallmatrix{\cdots & \alpha & \cdots & \alpha & \cdots}\\
        &=0.
    \end{align*}

    当 \(|S|=1\), 设 \(S=\{i\}\), 由按行展开定理, 有
    \begin{align*}
        &\vsmallmatrix{\beta_S(1) & \beta_S(2) & \cdots & \beta_S(n)}
        =\vsmallmatrix{\gamma_1 & \cdots & \gamma_{i-1} & b_i\alpha & \gamma_{i+1} & \cdots & \gamma_n}\\
        ={}&
        \begin{vmatrix}
            c_1 & & & a_1b_i & & &\\
            & \ddots & & \vdots & & &\\
            & & c_{i-1} & a_{i-1}b_i & & &\\
            & & & a_ib_i & & &\\
            & & & a_{i+1}b_i & c_{i+1} & &\\
            & & & \vdots & & \ddots &\\
            & & & a_nb_i & & & c_n
        \end{vmatrix}
        =\frac{a_ib_i}{c_i}\cdot(-1)^{i+i}\prod_{j=1}^nc_j
        =\frac{a_ib_i}{c_i}\prod_{j=1}^nc_j.
    \end{align*}

    当 \(S=\varnothing\), 有
    \[
    \vsmallmatrix{\beta_S(1) & \beta_S(2) & \cdots & \beta_S(n)}=\prod_{j=1}^nc_j.
    \]

    综上, 有
    \begin{align*}
        \det(a_{ij})
        &=\vsmallmatrix{b_1\alpha+\gamma_1 & b_2\alpha+\gamma_2 & \cdots & b_n\alpha+\gamma_n}\\
        &=\sum_{S\in\{1,2,\cdots,n\}}\vsmallmatrix{\beta_S(1) & \beta_S(2) & \cdots & \beta_S(n)}\\
        &=\prod_{j=1}^nc_j+\sum_{i=1}^n\frac{a_ib_i}{c_i}\prod_{j=1}^nc_j.\qedhere
    \end{align*}
\end{solution}